\scanpage{11}
- \mspage{8} -

où $\delta$ est le discriminant du second membre de (3.1). Le discriminant $\delta$ étant invariant par translation, il suffit de vérifier cette identité algébrique pour 3 inversible et $a_2=0$.

Lorsque $c_4$ est inversible, un et un seul choix de $x$ et $y$ donne $d_6=0$, i.e. $a_2a_4-9a_6=0$. On a, lorsque cette condition est remplie,
\[
(3.3)\qquad -c_6=2^6(a_2^3-3^3a_6)\qquad\Delta=2^6(-a_2^3a_6-a_4^3+2\cdot3^3a_6^2)
\]
et
\[
\begin{aligned}
(3.4)\qquad a_2c_4&=-2^{-2}c_6\\
a_4c_4^2&=-2^2\,3^2\,\Delta\\
a_6c_4^3&=c_6\Delta.
\end{aligned}
\]
Il suffit de vérifier ces identités en caractéristique 0, auquel cas la troisième est le produit des deux premières, puisque $d_6=0$.

\uline{Proposition} 3.5.\quad \uline{Au dessus de} $\mathrm{Spec}(\mathbb Z[2^{-1}])$, \uline{il existe un schéma de module pour les courbes de genre un munies d'une forme différentielle invariante inversible et telles que} $c_4$ \uline{soit inversible. Ce schéma est}
\[
\mathrm{Spec}(\mathbb Z[2^{-1}][a_2,a_4,a_6,(a_2^2-3a_4)^{-1}]/(a_2a_4-9a_6))
\]
\uline{avec pour courbe universelle} (\uline{en coordonnées non homogène}) \inserted{\unsure{$\rho$}}
\margin{section $e$}
\[
y^2=x^3+a_2x^2+a_4x+a_6
\]
\margin{10}
